\section*{Capítulo \ref{ch:b3:partition-functions} --- Termodinámica estadística: funciones de partición}

\begin{solution}{exo:b3:partition-functions:1}
Se escribe cada distribución como los cuantos que tienen los cuatro osciladores, de mayor
a menor. ``4, 0, 0, 0'': $4!/(3!\,1!) = 4$; ``3, 1, 0, 0'': $4!/(2!\,1!\,1!) = 12$; ``2, 2, 0, 0'':
$4!/(2!\,2!) = 6$; ``2, 1, 1, 0'': $4!/(1!\,2!\,1!) = 12$; ``1, 1, 1, 1'': 1. Total: 35 $= \binom73$.
Las más probables son ``3, 1, 0, 0'' y ``2, 1, 1, 0'' (12 cada una); la segunda, con
\omterm{def:b3:partition-functions:boltzmann}{poblaciones} 1, 2, 1 para 0, 1, 2 cuantos, ya disminuye con la energía como lo hace la
\omterm{def:b3:partition-functions:boltzmann}{distribución de Boltzmann} para números grandes.
\end{solution}

\begin{solution}{exo:b3:partition-functions:2}
$\eu^{-2500/(8.314 \times 298.15)} = \eu^{-1.009} = 0.365$; a \qty{1000}{K}, $\eu^{-0.301} = 0.740$. A temperatura
muy alta, la razón tiende a 1 (\omterm{def:b3:partition-functions:boltzmann}{poblaciones} iguales, nunca una inversión).
\end{solution}

\begin{solution}{exo:b3:partition-functions:3}
$m = 39.95\times10^{-3}/6.022\times10^{23} = \qty{6.634e-26}{kg}$; $\Lambda = h/\sqrt{2\pi mkT} =
\qty{1.600e-11}{m} = \qty{16.0}{pm}$; $q^{\mathrm T} = 10^{-3}/(1.600\times10^{-11})^3 = \num{2.44e29}$.
\end{solution}

\begin{solution}{exo:b3:partition-functions:4}
$hc/k = \qty{1.4388}{cm.K}$. \ce{N2}: $\theta_{\mathrm R} = \qty{2.863}{K}$, $q^{\mathrm R} = 298.15/(2 \times 2.863) = 52.1$.
\ce{HCl}: $\theta_{\mathrm R} = \qty{15.02}{K}$, $q^{\mathrm R} = 298.15/15.02 = 19.85$ (suma exacta: 20.19).
\end{solution}

\begin{solution}{exo:b3:partition-functions:5}
\ce{I2}: $\theta_{\mathrm V} = 1.4388 \times 213.27 = \qty{306.9}{K}$, $q^{\mathrm V} = 1/(1 - \eu^{-1.029}) = 1.556$;
fracción excitada: $1 - 1/q^{\mathrm V} = 0.357$. \ce{N2}: $\theta_{\mathrm V} = \qty{3352}{K}$, $q^{\mathrm V} = 1.000013$,
fracción excitada: \num{1.3e-5}.
\end{solution}

\begin{solution}{exo:b3:partition-functions:6}
\ce{NO}: $q^{\mathrm E} = 2 + 2\eu^{-119.73/207.22} = 2 + 2 \times 0.561 = 3.12$; fracción superior: $1.12/3.12
= 0.36$. A alta temperatura, $q^{\mathrm E} \to 4$ (dos niveles igualmente poblados). \ce{Cl}:
$q^{\mathrm E} = 4 + 2\eu^{-882.35/207.22} = 4 + 2 \times 0.0142 = 4.03$.
\end{solution}

\begin{solution}{exo:b3:partition-functions:7}
$\ln q = \ln V + \frac32\ln T + \text{cte}$, luego $U - U(0) = NkT^2 \times \frac{3}{2T} = \frac32NkT$; para un
mol, $\frac32RT = \qty{3.72}{kJ/mol}$ a \qty{298.15}{K}, y $C_V = \frac32R = \qty{12.47}{J.K^{-1}.mol^{-1}}$.
\end{solution}

\begin{solution}{exo:b3:partition-functions:8}
$kT/p^\circ = \qty{4.116e-26}{m^3}$, $\Lambda^3 = \qty{4.093e-33}{m^3}$, razón \num{1.006e7}:
$S_{\mathrm m} = 8.3145 \times (16.124 + 2.5) = \qty{154.85}{J.K^{-1}.mol^{-1}}$, el valor tabulado hasta la
última cifra. Dividir $p$ por 10 añade $R\ln10 = \qty{19.14}{J.K^{-1}.mol^{-1}}$.
\end{solution}

\begin{solution}{exo:b3:partition-functions:9}
$S^{\mathrm R}_{\mathrm m} = R[\ln(298.15/15.02) + 1] = 8.3145 \times (2.988 + 1) = \qty{33.16}{J.K^{-1}.mol^{-1}}$.
\end{solution}

\begin{solution}{exo:b3:partition-functions:10}
$\frac{\dd}{\dd J}\bigl[(2J+1)\eu^{-\theta J(J+1)/T}\bigr] = \bigl[2 - (2J+1)^2\theta/T\bigr]\eu^{-\theta J(J+1)/T} = 0$ da
$2J + 1 = \sqrt{2T/\theta}$, es decir, $J_{\max} = \sqrt{T/2\theta} - \frac12$. \ce{HCl}: 2.66, luego $J = 3$; \ce{CO}:
6.86, luego $J = 7$. La intensidad de una raya de la \omterm{def:b3:rovibrational-spectroscopy:branches}{rama P} o R sigue la
\omterm{def:b3:partition-functions:boltzmann}{población} de su nivel inferior, que es máxima en $J_{\max}$, a unas pocas rayas del hueco
del centro de la banda.
\end{solution}

\begin{solution}{exo:b3:partition-functions:11}
$\int_0^\infty f\dd J = T/\theta$, $f(0) = 1$, $f'(0) = 2 - \theta/T$. Por tanto, $q^{\mathrm R} \approx T/\theta + \frac12 -
\frac16 + \frac{\theta}{12T} \approx T/\theta + \frac13$. El error relativo de $T/\theta$ es de alrededor de $\theta/3T$:
inferior al 1\,\% cuando $T > 33\,\theta$. Para \ce{H2}, $\theta = 1.4388 \times 59.32 = \qty{85.3}{K}$: a
\qty{298.15}{K} el error es del 10\,\%, no aceptable (y la restricción sobre $J$ debida al espín nuclear
también debe tratarse exactamente, \cref{ch:b3:statistical-thermo-applied}).
\end{solution}

\begin{solution}{exo:b3:partition-functions:12}
Para $x = \theta_{\mathrm V}/T \to \infty$, $x/(\eu^x - 1) \to 0$ y $\ln(1 - \eu^{-x}) \to 0$: $S \to 0$. Para $x \to 0$,
$x/(\eu^x - 1) \approx 1 - x/2$ y $-\ln(1 - \eu^{-x}) \approx -\ln x + x/2$: $S \to R(1 - \ln x) = R[1 + \ln(T/\theta_{\mathrm V})]$.
\ce{I2}, $x = 1.029$: valor exacto $R(0.5721 + 0.4420) = \qty{8.43}{J.K^{-1}.mol^{-1}}$; límite $R(1 - 0.0289) =
\qty{8.07}{J.K^{-1}.mol^{-1}}$. Incluso con $T \approx \theta_{\mathrm V}$, el límite clásico solo se queda un 4\,\% por debajo.
\end{solution}

\begin{solution}{pb:b3:partition-functions:1}
\textbf{1.} $m = 28.014\times10^{-3}/6.02214\times10^{23} = \qty{4.652e-26}{kg}$.
\textbf{2.} $\Lambda = h/\sqrt{2\pi mkT} = \qty{1.910e-11}{m}$ (\qty{19.1}{pm}).
\textbf{3.} $kT/p^\circ = 1.380649\times10^{-23} \times 298.15/10^5 = \qty{4.116e-26}{m^3}$.
\textbf{4.} $q^{\mathrm T}/N = 4.116\times10^{-26}/(1.910\times10^{-11})^3 = \num{5.905e6}$: unos seis millones
de estados de traslación por molécula, de modo que dos moléculas casi nunca comparten uno y
$Q = q^N/N!$ es válida.
\textbf{5.} $S^{\mathrm T}_{\mathrm m} = R(\ln\num{5.905e6} + \frac52) = 8.3145 \times 18.091 = \qty{150.42}{J.K^{-1}.mol^{-1}}$.
\textbf{6.} $+R\ln2 = \qty{5.76}{J.K^{-1}.mol^{-1}}$, es decir, \num{156.18}.
\textbf{7.} $B_0 = 1.998241 - 0.008659 = \qty{1.989582}{cm^{-1}}$.
\textbf{8.} $\theta_{\mathrm R} = 1.438777 \times 1.989582 = \qty{2.8626}{K}$.
\textbf{9.} $\sigma = 2$: los dos núcleos son idénticos; dar la vuelta a la molécula de extremo a extremo
da una orientación indistinguible, y la simetría de intercambio de los
núcleos solo permite, para cada estado de espín nuclear, la mitad de los niveles de rotación.
\textbf{10.} $q^{\mathrm R} = 298.15/(2 \times 2.8626) = 52.08$.
\textbf{11.} Sí: $T/\theta_{\mathrm R} = 104$, un error de alrededor del 0.3\,\% en $q^{\mathrm R}$ y mucho menor en la
entropía.
\textbf{12.} $S^{\mathrm R}_{\mathrm m} = R(\ln52.08 + 1) = 8.3145 \times 4.9527 = \qty{41.18}{J.K^{-1}.mol^{-1}}$.
\textbf{13.} $J_{\max} = \sqrt{298.15/5.725} - 0.5 = 6.72$: $J = 7$ (\omterm{def:b3:partition-functions:boltzmann}{población} 0.0839, justo por encima de
0.0831 para $J = 6$).
\textbf{14.} $2358.57 - 2 \times 14.324 = \qty{2329.92}{cm^{-1}}$.
\textbf{15.} $\theta_{\mathrm V} = 1.438777 \times 2329.92 = \qty{3352}{K}$.
\textbf{16.} $x = 3352.2/298.15 = 11.24$; $q^{\mathrm V} = 1/(1 - \eu^{-11.24}) = 1.000013$.
\textbf{17.} $\eu^{-11.24}/q^{\mathrm V} = \num{1.3e-5}$.
\textbf{18.} $S^{\mathrm V}_{\mathrm m} = R[11.24\eu^{-11.24} + \eu^{-11.24}] = \qty{0.0013}{J.K^{-1}.mol^{-1}}$, despreciable.
\textbf{19.} $g_0 = 1$ (${}^1\Sigma$): $R\ln1 = 0$. Los estados electrónicos excitados están varios
electronvoltios más arriba, más de cien veces $kT$: sus factores de Boltzmann son
completamente despreciables.
\textbf{20.} $150.42 + 41.18 + 0.00 + 0 = \qty{191.60}{J.K^{-1}.mol^{-1}}$.
\textbf{21.} La diferencia es de \qty{0.01}{J.K^{-1}.mol^{-1}}, del orden de las
aproximaciones hechas (\omterm{def:b3:quantum-model-systems:rotor}{rotor rígido}, separación de los movimientos).
\textbf{22.} Las capacidades caloríficas del sólido medidas desde unos pocos kelvin (extrapoladas
como $T^3$ por debajo), las entalpías de la transición sólido--sólido, de fusión y
de vaporización divididas por sus temperaturas, la capacidad calorífica del gas,
y una corrección por la no idealidad.
\textbf{23.} Los núcleos se conservan en toda reacción, de modo que su entropía de espín es la
misma a ambos lados y se cancela; la calorimetría tampoco la ve, ya que los
espines nucleares siguen desordenados en el cristal hasta las temperaturas más bajas
medidas.
\textbf{24.} \textbf{$S^\circ(\ce{N2}, \qty{298.15}{K}) = \qty{191.6}{J.K^{-1}.mol^{-1}}$ a partir de las constantes espectroscópicas},
de los cuales 150.4 proceden de la traslación y 41.2 de la rotación.
\end{solution}
